\lyseis{7.8}{Δυνάμεις με εκθέτη φυσικό}
\sol{1} α) $32$ \quad β) $81$ \quad γ) $-8$ \quad δ) $16$ \quad ε) $-1$ \quad στ) $0$.
\sol{2} α) $9$ \quad β) $-9$ \quad γ) $-\frac{1}{8}$ \quad δ) $0{,}001$ \quad ε) $0{,}04$ \quad στ) $\frac{8}{27}$.
\sol{3} α) Θετικό (άρτιος εκθέτης). β) Αρνητικό. γ) Αρνητικό ($(-3)^4=81$ και μετά το μείον). δ) Αρνητικό ($(-2)^3=-8$ και $(-3)^2=9$).
\sol{4} α) $2^7$ \quad β) $5^3$ \quad γ) $3^6$ \quad δ) $7^6$ \quad ε) $(-2)^7$ \quad στ) $(-4)^2=16$.
\sol{5} α) $10^3=1\,000$ \quad β) $(-6)^2=36$ \quad γ) $\frac{8}{27}$ \quad δ) $\frac{9}{16}$.
\sol{6} α) $18+8-2=24$ \quad β) $12-16:(-2)=12+8=20$ \quad γ) $5-12=-7$ \quad δ) $(-1+1)\cdot10=0$.
\sol{7} α) Λάθος: $8$. β) Σωστό. γ) Λάθος: $-16$. δ) Σωστό. ε) Λάθος: $\alpha^5$. στ) Σωστό. ζ) Λάθος.
\sol{8} Πλήθη: $1,\,2,\,4,\,8,\,16,\,32,\,64,\,128,\,256,\,512,\,1\,024$. Ξεπερνά τα $1\,000$ μετά από $10$ ώρες.
\sol{9} Όγκος $12^3=1\,728$ cm$^3$. Επιφάνεια $6\cdot144=864$ cm$^2$.
\sol{10} Στο $8$ο: $2^7=128$. Στο $11$ο: $2^{10}=1\,024$.
\sol{11} $3^5=243$ βακτήρια.
\sol{12} Ο $20$ός είναι $20^2=400$ και ο $\nu$-οστός $\nu^2$.
\sol{13} $3\cdot2^4=48$ €.
\sol{14} $10^3\cdot10^4=10^7$ και $10^8:10^5=10^3$.
\sol{15} Το $2\cdot3$ είναι $2+2+2=6$ (πρόσθεση του $2$ τρεις φορές). Το $2^3=2\cdot2\cdot2=8$ (πολλαπλασιασμός του $2$ τρεις φορές).
\sol{16} $(3+4)^2=49$ ενώ $3^2+4^2=9+16=25$.
\sol{17} Όταν ο εκθέτης $\nu$ είναι άρτιος $(-\alpha)^\nu=\alpha^\nu$, όταν είναι περιττός $(-\alpha)^\nu=-\alpha^\nu$.
\sol{18} $1^\nu=1$, $0^\nu=0$ και $(-1)^\nu=1$ για άρτιο και $-1$ για περιττό $\nu$.
\sol{19} $2^3\cdot2^4=(2\cdot2\cdot2)\cdot(2\cdot2\cdot2\cdot2)=2\cdot2\cdot2\cdot2\cdot2\cdot2\cdot2=2^7$.
\sol{20} $3^4=81$ και $4^3=64$, άρα το $3^4$.
\sol{21} $2,\,4,\,8,\,16,\,32$ στρώματα, δηλαδή $2^\nu$.
